\chapter{Marco teórico}

\section{Algoritmos de ordenamiento}

Existen muchísimos algoritmos de ordenamiento, desde algunos muy eficientes, que pueden llegar a ser de una complejidad temporal de apenas $O(n)$, 
como \textit{Radix Sort}, hasta algunos extremadamente ineficientes, que pueden llegar incluso a ser de $O(n \times n!)$, como lo es el caso de 
\textit{Bogo Sort}, pero en este trabajo nos centraremos principalmente en \textit{Merge Sort} y \textit{Quick Sort}, que son algoritmos relativamente
eficientes y que emplean el principio de \textbf{"Divide y Vencerás"}. \\

\subsection{Merge Sort}

\textit{Merge Sort}, también conocido como "Ordenamiento por Mezcla", es un algoritmo de ordenamiento recursivo, que en cada paso va dividiendo el 
arreglo en 2 subarreglos de aproximadamente la mitad del tamaño anterior, y eso lo hace hasta que todos los subarreglos sean de tamaño 1. Cuando eso sucede,
se comienzan a combinar pares de subarreglos adyacentes, recorriéndolos linealmente para fusionarlos en el orden correcto, hasta que finalmente todo el 
arreglo queda ordenado. \\

En base al principio en el que se fundamenta \textit{Merge Sort}, es posible deducir su complejidad temporal mediante la siguiente ecuación de recurrencia: \\

$T(n) = 2 T(\frac{n}{2}) + O(n)$ \\

Esto se debe a que el problema se divide recursivamente en 2 subproblemas de tamaño $n/2$, mientras que el costo adicional corresponde al proceso de 
fusión de los subarreglos, el cual requiere recorrer linealmente los elementos para ubicarlos en el orden correcto, por lo que dicho costo es de 
complejidad lineal. \\

Para resolver la ecuación de recurrencia, se puede emplear el método de iteración. Considerando una constante $k$ para representar el costo lineal, se tiene: \\

$T(n) = 2 T(\frac{n}{2}) + kn$ \\

A partir de esto, la expresión equivalente para $T(\frac{n}{2})$ sería: \\

$T(\frac{n}{2}) = 2 T(\frac{n}{4}) + k(\frac{n}{2})$ \\

Entonces, reemplazando en la ecuación original, se obtiene: \\

$T(n) = 2 \left( 2 T(\frac{n}{4}) + k(\frac{n}{2}) \right) + kn$ \\

Simplificando: \\

$T(n) = 4T(\frac{n}{4}) + 2kn$ \\

Ahora, realizando otra iteración, se tiene que: \\

$T(\frac{n}{4}) = 2 T(\frac{n}{8}) + k(\frac{n}{4})$ \\

Reemplazando nuevamente: \\

$T(n) = 4 \left( 2T(\frac{n}{8}) + k(\frac{n}{4}) \right) + 2kn$ \\

$T(n) = 8T(\frac{n}{8}) + 3kn$ \\

Continuando de esta forma, se puede generalizar que: \\

$T(n) = 2T(\frac{n}{2}) + kn = 4T(\frac{n}{4}) + 2kn = 8T(\frac{n}{8}) + 3kn = 2^i T(\frac{n}{2^i}) + ikn$ \\

La recursión finaliza cuando se alcanza el caso base, es decir, cuando los subarreglos tienen tamaño 1. Por lo tanto, debe cumplirse que: \\

$\frac{n}{2^i} = 1$ \\

Despejando: \\

$i = \log_{2}(n)$ \\

Reemplazando esto en la expresión general obtenida anteriormente: \\

$T(n) = 2^{\log_{2}(n)} T(1) + \log_{2}(n)kn$ \\

Como $2^{\log_{2}(n)} = n$, entonces: \\

$T(n) = nT(1) + kn\log_{2}(n)$ \\

Finalmente, considerando que $T(1)$ y $k$ son constantes, se concluye que: \\

$T(n) \in O(n \log n)$ \\

Por lo tanto, \textit{Merge Sort} posee complejidad temporal $O(n \log n)$ tanto en el mejor caso, como en el caso promedio y peor caso, ya que el arreglo 
siempre se divide de manera uniforme, independientemente del orden inicial de los elementos. \\

\subsection{Merge-Insertion Sort}

\textit{Merge-Insertion Sort} corresponde a una variante optimizada de \textit{Merge Sort}, cuyo objetivo es reducir el overhead recursivo que se produce al 
trabajar con subarreglos muy pequeños. \\

La idea principal consiste en que, cuando los subarreglos alcanzan un tamaño suficientemente pequeño, en lugar de continuar dividiéndolos recursivamente, 
se utiliza \textit{Insertion Sort} para ordenarlos directamente. \\

Esto se hace debido a que \textit{Insertion Sort}, a pesar de poseer complejidad $O(n^2)$ en el peor caso, tiene un rendimiento muy eficiente para arreglos 
de tamaño reducido, ya que posee un overhead muy bajo y además aprovecha favorablemente arreglos parcialmente ordenados. \\

De esta manera, \textit{Merge-Insertion Sort} busca combinar las ventajas de ambos algoritmos: la eficiencia asintótica de \textit{Merge Sort} para arreglos 
grandes, junto con la simplicidad y rapidez de \textit{Insertion Sort} en subarreglos pequeños. \\

En términos de complejidad asintótica, esta variante sigue perteneciendo a $O(n \log n)$, aunque en la práctica puede presentar mejoras de rendimiento 
respecto a la versión clásica de \textit{Merge Sort}, dependiendo del umbral utilizado para decidir cuándo aplicar \textit{Insertion Sort}, que es uno de
los aspectos que se analizará en el presente trabajo. \\

\subsection{Quick Sort}

